% TOP_PROBLEM: {"id": 2183, "title": "Erdős Problem #567", "category_id": 3, "category": {"id": 3, "name": "graph_theory", "display_name": "Graph Theory", "description": "Problems involving graphs, networks, and their properties.", "slug": "graph-theory", "order_index": 3, "created_at": "2026-07-31T15:26:25.671Z"}, "status": "open"}
% FIRST_POSTED: null
% ENTRY_KIND: statement_only
% Imported from: unsolvedmath_additional_500.tex; UnsolvedMath ID 2183
% !TeX program = lualatex
% Compile twice: lualatex 2183.tex
% Self-contained: no external bibliography, images, or \input files.
% Numeric section labels and visible numbers are original UnsolvedMath IDs.
\documentclass[11pt,a4paper]{article}
\usepackage[margin=24mm,headheight=14pt]{geometry}
\usepackage{fontspec}
\setmainfont{Latin Modern Roman}
\setsansfont{Latin Modern Sans}
\setmonofont{Latin Modern Mono}
\usepackage{amsmath,amssymb,mathtools}
\usepackage{unicode-math}
\setmathfont{Latin Modern Math}
\usepackage{microtype}
\usepackage{enumitem,xurl,needspace}
\usepackage[dvipsnames]{xcolor}
\usepackage{hyperref}
\hypersetup{unicode=true,colorlinks=true,linkcolor=MidnightBlue,urlcolor=MidnightBlue,
 pdftitle={UnsolvedMath Problem 2183},pdfauthor={Source-annotated compilation},bookmarksnumbered=true}
\usepackage{bookmark,tocloft,titlesec,fancyhdr}
\setlength{\cftsecnumwidth}{4.2em}
\cftsetpnumwidth{2.5em}
\cftsetrmarg{4em}
\titleformat{\section}{\Large\bfseries\raggedright}{\thesection}{0.7em}{}
\pagestyle{fancy}\fancyhf{}
\fancyhead[L]{\small\sffamily Additional UnsolvedMath statements}
\fancyhead[R]{\small Original catalogue IDs}
\fancyfoot[C]{\thepage}
\setlength{\parindent}{0pt}
\setlength{\parskip}{5pt plus 1pt}
\setlength{\emergencystretch}{3em}
\setcounter{tocdepth}{1}\setcounter{secnumdepth}{1}
\setlist[itemize]{leftmargin=3em,itemsep=4pt,topsep=4pt,parsep=0pt}
\allowdisplaybreaks
\newcommand{\N}{\mathbb{N}}
\newcommand{\Z}{\mathbb{Z}}
\newcommand{\Q}{\mathbb{Q}}
\newcommand{\R}{\mathbb{R}}
\newcommand{\C}{\mathbb{C}}
\newcommand{\F}{\mathbb{F}}
\newcommand{\A}{\mathbb{A}}
\newcommand{\PP}{\mathbb{P}}
\newcommand{\E}{\mathbb{E}}
\newcommand{\eps}{\varepsilon}
\newcommand{\dd}{\,\mathrm d}
\newcommand{\sref}[2]{\hyperlink{source:#1:#2}{[S#2]}}
\newif\ifshowcatalogue
\showcataloguetrue
\newcommand{\cataloguescope}[1]{\ifshowcatalogue
 \paragraph{Statement recorded in the supplied source.}
 #1\fi}
\begin{document}
% UnsolvedMath ID 2183; source catalogue code EP-567

\setcounter{section}{2182}

\section{Linear Off-Diagonal Ramsey Bounds for Three Fixed Graphs}\label{2183}

{\small\sffamily UnsolvedMath ID 2183 · Graph Theory}

\subsection{Definitions and mathematical statement}

\paragraph{Shared notation.}Unless a section says otherwise, $\N=\{1,2,\ldots\}$,
$\Z$ is the ring of integers, $\Q,\R,\C$ are the rational, real, and complex
fields, $[N]=\{1,\ldots,\lfloor N\rfloor\}$, and $\log$ is the natural
logarithm. For real functions of a parameter tending to infinity,
$f\ll g$ and $f=O(g)$ mean $|f|\leq Cg$ eventually for some fixed $C$;
$f\gg g$ reverses this comparison; $f=o(g)$ means $f/g\to0$;
$f\sim g$ means $f/g\to1$; and $f\asymp g$ means both order bounds.
A subscript on $O$ or $\ll$ specifies allowed constant dependence.
The expression $n^{o(1)}$ in an upper bound means growth smaller than
$n^\epsilon$ for each fixed $\epsilon>0$. Local definitions govern any
exception, finite-field convention, or use of the same symbol for another
object. An asymptotic claim is not automatically an effective algorithm.



The graph $H_5$ has the edges of a five-cycle together with two chords having no common endpoint. The other fixed graphs are the three-cube and the complete bipartite graph with parts of size three. The constant may depend on the chosen fixed graph, but not on $H$.

A graph has a vertex set $V$ and edges that are unordered pairs of distinct vertices. A subgraph need not be induced unless that word is stated. A proper vertex colouring gives adjacent vertices different colours; $\chi(G)$ is the least number of colours.

$R(G,H)$ is the least order of a complete graph whose every red--blue edge colouring has a red copy of $G$ or a blue copy of $H$. For integer arguments, $R(s,t)=R(K_s,K_t)$. Write $R(G)=R(G,G)$ and $R(k)=R(K_k,K_k)$. The notation $R(G;q)$ instead uses $q$ colours and the same target in each colour.

$K_t$ is a complete graph and $C_t$ a simple cycle of length $t$. A complete multipartite graph has no edges within parts and all edges between different parts. An independent set has no internal edge; a clique has all its possible internal edges. \sref{2183}{1} \sref{2183}{2}

\paragraph{Statement recorded in the supplied source.}
Let $G$ be either $Q_3$ or $K_{3,3}$ or $H_5$ (the last formed by adding two vertex-disjoint chords to $C_5$). Is it true that, if $H$ has $m$ edges and no isolated vertices, then $ R(G,H)\ll m? $ \sref{2183}{1}

\paragraph{Residual task reported by the archive.}

The embedded review (2026-08-17) records: Prove or refute Ramsey size-linearity for each of Q\_3, K\_{3,3}, and H\_5 against arbitrary graphs without isolated vertices. \sref{2183}{1}

{\small\emph{Transcription note.} Only unambiguous escape/formatting damage and nonmathematical serialization debris were repaired in the displayed source statement.}

\subsection{Short English statement}

For each of the three specified small graphs, does a constant times the opposing graph's edge count suffice in the off-diagonal Ramsey problem?

\subsection{Sources}

{\small

\begin{itemize}

\item[\hypertarget{source:2183:1}{S1}] ULAM.AI, \emph{UnsolvedMath}, supplied \texttt{problems.json}, record 2183 (EP-567), statement and background. The embedded literature-review date is 2026-08-17; that is the archive’s review date, not a fresh certification. Dataset landing page: \url{https://huggingface.co/datasets/ulamai/UnsolvedMath}.

\item[\hypertarget{source:2183:2}{S2}] T. F. Bloom, \emph{Erdős Problems}, Problem 567, \url{https://www.erdosproblems.com/567}. Reference imported from the supplied record; the live problem page was not independently checked for this compilation.

\item[\hypertarget{source:2183:3}{S3}] Domagoj Bradač, Lior Gishboliner, and Benny Sudakov, On Ramsey Size-Linear Graphs and Related Questions, SIAM J. Discrete Math. 38 (2024), 225-242, \nolinkurl{doi:10.1137/22M1481713.} \url{https://arxiv.org/abs/2202.10388}. Bibliographic information imported from the supplied record.

\item[\hypertarget{source:2183:4}{S4}] Brada\'C, D. and Gishboliner, L. and Sudakov, B., On Ramsey size-linear graphs and related questions. arXiv:2202.10388 (2023). Bibliographic information imported from the supplied record.

\item[\hypertarget{source:2183:5}{S5}] Erd\H{o}s, Paul, Some of my favourite problems in number theory, combinatorics, and geometry. Resenhas (1995), 165-186. Bibliographic information imported from the supplied record.

\end{itemize}

}

\label{end:2183}
\end{document}
